\chapter{Semantic Manifold Foundations}
\label{ch:semantic-manifold}

This chapter establishes the geometric, categorical, and informational
foundations for the RSVP--Polyxan unified hyperstructure. It introduces the
\emph{semantic manifold} on which the RSVP fields $(\Phi,\mathbf{v},S)$ live,
defines its smooth and stratified structures, constructs the jet bundles needed
for action principles, and shows how Polyxan hypermedia atoms embed as sheaf
sections over semantic open sets.

The semantic manifold $\mathcal{M}$ is the base space for:
\begin{itemize}[leftmargin=2em]
  \item differential geometric RSVP field theory (Parts I and III),
  \item Polyxan typed hypergraph constructions (Part II),
  \item semantic embeddings (Parts III and VII),
  \item formal verification (Part IV),
  \item system engineering and simulation architecture (Parts V and VI).
\end{itemize}

The goal of this chapter is to provide a rigorous mathematical substrate on
which the remainder of the monograph can rest. Later constructions will rely
on the structures defined here, especially the stratification $\mathcal{S}$,
the semantic sheaf $\mathscr{S}$, and the semantic metric~$g$.

% ================================================================
\section{The Semantic Base Manifold}

We begin by defining the space in which semantic objects ``live''.

\begin{definition}[Semantic Manifold]
A \emph{semantic manifold} is a paracompact, second-countable differentiable
manifold $\mathcal{M}$ equipped with:
\begin{enumerate}
  \item a smooth atlas $\mathcal{A} = \{(U_i,\varphi_i)\}$;
  \item a (possibly indefinite) semantic metric $g$;
  \item a stratification $\mathcal{S}$ into semantic layers;
  \item a semantic sheaf $\mathscr{S}$ whose sections encode local semantic
        structures, Polyxan atoms, and RSVP configurations.
\end{enumerate}
\end{definition}

Unlike physical spacetime, points of $\mathcal{M}$ represent semantic
situations. Coordinates are semantic, not spatial.

\subsection{Motivation: Why a Manifold?}

Semantic content often exhibits:
\begin{itemize}
  \item local linearity (e.g., embedding models);
  \item global curvature (e.g., clustering of concepts);
  \item local triviality (via charts);
  \item smooth transitions of meaning (via geodesic flows).
\end{itemize}

The manifold formalism captures these properties naturally and supports the
variational principles needed for RSVP.

% ================================================================
\section{Charts and Semantic Coordinates}

Each chart $(U,\varphi)$ introduces local coordinates
$ (x^1, \dots, x^n) $, but meaningful coordinates arise from RSVP’s canonical
decomposition:

\begin{definition}[Semantic Coordinate System]
A chart $(U,\varphi)$ admits a canonical coordinate decomposition
\[
(x^\mathrm{aff}, x^\mathrm{geom},
  x^\mathrm{dyn}, x^\mathrm{poly})
\]
corresponding to:
\begin{itemize}
  \item \textbf{affective coordinates} $x^\mathrm{aff}$ — valence, salience,
        motivational gradients;
  \item \textbf{geometric coordinates} $x^\mathrm{geom}$ — structural or
        embedding space directions;
  \item \textbf{dynamical coordinates} $x^\mathrm{dyn}$ — transition or flow
        parameters related to $\mathbf{v}$;
  \item \textbf{Polyxan coordinates} $x^\mathrm{poly}$ — local indices for
        hypergraph embeddings and content atoms.
\end{itemize}
\]
\end{definition}

This decomposition embeds RSVP field structure directly into the local
coordinate system, enabling coherent coupling between manifold fields and
hypergraph structures.

\subsection{Change of Coordinates and Semantic Jacobians}

Coordinate transformations
\[
\psi = \varphi_j \circ \varphi_i^{-1} : \varphi_i(U_i \cap U_j)
\to \varphi_j(U_i \cap U_j)
\]
must preserve the semantic decomposition.

\begin{proposition}
If semantic coordinates are preserved under chart transitions, then the Jacobian
matrix of $\psi$ is block-diagonal up to permutation.
\end{proposition}

This ensures that affective, geometric, dynamical, and Polyxan coordinates
transform compatibly, which is crucial for coupling RSVP with Polyxan indices.

% ================================================================
\section{Jet Extensions and Higher-Order Structures}

The RSVP Lagrangian requires second-order derivatives of $\Phi$ and first-order
derivatives of $\mathbf{v}$ and~$S$.

\begin{definition}[Jet Bundle]
The $k$-jet bundle $J^k(\mathcal{M},\mathbb{R})$ is the fiber bundle whose fiber
at $x$ consists of equivalence classes of smooth functions agreeing to
$k$th order at~$x$.
\end{definition}

Relevant bundles include:
\[
J^2(\Phi),
\qquad
J^1(\mathbf{v}),
\qquad
J^1(S).
\]

These provide the natural domain for the RSVP action functional (Ch.~2).

\subsection{Jet Compatibility with Stratification}

\begin{lemma}[Jet--Stratification Compatibility]
If $\mathcal{M}$ is Whitney stratified into smooth strata $\{X_\alpha\}$, then
each $J^k(\mathcal{M})$ restricts to a smooth bundle on each $X_\alpha$, and
restriction maps commute with jet prolongations.
\end{lemma}

\begin{proof}
Since each stratum $X_\alpha$ is a smooth manifold, jets restrict smoothly.
Whitney conditions ensure transverse smoothness between strata; see standard
results in Thom--Mather theory.
\end{proof}

This lemma ensures that RSVP variation principles remain well-posed even when
$\mathcal{M}$ exhibits abrupt semantic discontinuities across strata.

% ================================================================
\section{RSVP Fields on the Semantic Manifold}

The RSVP framework introduces three primary fields:
\begin{align*}
\Phi &: \mathcal{M} \to \mathbb{R},
&& \text{scalar semantic potential}, \\
\mathbf{v} &: \mathcal{M} \to T\mathcal{M},
&& \text{semantic drift / agency field}, \\
S &: \mathcal{M} \to \mathbb{R}_{\ge 0},
&& \text{entropy density / uncertainty}.
\end{align*}

These are collected into a configuration tuple $\Psi = (\Phi,\mathbf{v},S)$.

\subsection{Semantic Metric and Covariant Derivatives}

The semantic metric $g$ measures the ``cost of semantic displacement.'' It may
be indefinite or degenerate on lower strata.

\begin{definition}[Semantic Gradient]
The gradient of $\Phi$ with respect to~$g$ is
\[
\nabla\!\Phi = g^{ij}
\,\partial_j \Phi\, \partial_i .
\]
\end{definition}

\begin{definition}[Semantic Laplacian]
The Laplacian of a scalar field is
\[
\Delta\Phi := \nabla_i \nabla^i \Phi.
\]
\end{definition}

These operators will appear in the RSVP Euler--Lagrange equations (Ch.~3).

% ================================================================
\section{Stratification and Semantic Layers}

Semantic content is rarely uniform across $\mathcal{M}$; it clusters in
regions of varying granularity.

\begin{definition}[Semantic Stratification]
A \emph{semantic stratification} of $\mathcal{M}$ is a filtration
\[
\mathcal{M} = X_0 \supset X_1 \supset X_2 \supset \cdots
\]
where each $X_i$ is a closed subset, and the strata
$S_i = X_i \smallsetminus X_{i+1}$
are smooth submanifolds.
\end{definition}

Interpretation:
\[
S_0 = \text{coarse semantic regions}, \qquad
S_1 = \text{topics or domains}, \qquad
S_2 = \text{concept clusters}, \qquad
S_3 = \text{atomic content}.
\]

This stratification underpins sheaf-theoretic reasoning (Ch.~10) and allows
Polyxan hypergraphs to adapt their ``resolution'' to the semantic layer.

\subsection{Stratified Metrics and Degenerate Directions}

The semantic metric $g$ may be degenerate along lower strata. Formally:

\begin{proposition}
If $g$ degenerates along $S_i$ but remains positive-definite on normal bundles
to $S_i$, then RSVP flows remain well-defined and transverse to strata.
\end{proposition}

This ensures that $\mathbf{v}$ and $\nabla \Phi$ need not be tangent to each
stratum, permitting cross-layer semantic jumps.

% ================================================================
\section{Sheaves of Semantic Structures}

The sheaf $\mathscr{S}$ organizes semantic information over $\mathcal{M}$.

\begin{definition}[Semantic Sheaf]
The \emph{semantic sheaf} $\mathscr{S}$ assigns to each open
$U \subseteq \mathcal{M}$ a category $\mathscr{S}(U)$ of semantic objects
(sections encoding meaning, embeddings, interpretations), with restriction
functors satisfying standard sheaf axioms.
\end{definition}

Unlike classical sheaves of sets, $\mathscr{S}$ is frequently:
\begin{itemize}
  \item a sheaf of categories,
  \item or a sheaf of $(\infty,1)$-categories (in later chapters),
  \item or a stack (in the Polyxan quine construction).
\end{itemize}

\subsection{Polyxan Atoms as Sections}

A Polyxan atom $a$ is a section
\[
a \in \Gamma(U,\mathscr{S}) = \mathscr{S}(U)
\]
defined over a region $U$ of semantic support.

\begin{definition}[Semantic Support]
The semantic support of a Polyxan atom is the minimal open set
\[
\operatorname{supp}(a)
\subseteq \mathcal{M}
\]
such that $a$ admits a section over it.
\end{definition}

This support determines where the RSVP fields couple to $a$ via curvature,
flows, and embedding dynamics.

% ================================================================
\section{Categorical Structures: Fibrations and Indexed Families}

To integrate Polyxan structures with RSVP dynamics, we introduce a categorical
framework.

\begin{definition}[Semantic Fibration]
A functor $\pi:\mathcal{E}\to\mathcal{M}$ is a
\emph{semantic fibration} if:
\begin{itemize}
  \item each fiber $\mathcal{E}_x$ is a category of semantic states at $x$;
  \item curves in $\mathcal{M}$ lift to functors between fibers;
  \item RSVP flows generate natural transformations between these functors.
\end{itemize}
\end{definition}

Polyxan hypergraphs appear as global sections of~$\pi$.

\subsection{Parallel Transport of Semantic Objects}

Given $\mathbf{v}$, parallel transport in $\pi$ gives a semantics-preserving
evolution of hypergraph nodes along RSVP flows:
\[
\mathsf{PT}_t : \mathcal{E}_{x_0} \to \mathcal{E}_{\varphi_t(x_0)}.
\]

This is the categorical analog of semantic drift modeled by $\mathbf{v}$.

% ================================================================
\section{Semantic Vector Fields and Flows}

\begin{definition}[Semantic Flow]
Given a smooth vector field $\mathbf{v} \in \Gamma(T\mathcal{M})$, its
associated flow is the one-parameter family of diffeomorphisms
$\{\varphi_t: \mathcal{M} \to \mathcal{M}\}$ satisfying
\[
\frac{d}{dt}\varphi_t(x) = \mathbf{v}(\varphi_t(x)),
\qquad \varphi_0 = \mathrm{id}.
\]
\end{definition}

\begin{proposition}[Completeness Criterion]
If $g(\mathbf{v},\mathbf{v})$ is bounded on compact sets, the
semantic flow is complete.
\end{proposition}

This ensures that semantic drift cannot ``blow up'' in finite time.

% ================================================================
\section{Semantic Curvature and Energy Bounds}

The curvature tensor of $g$ encodes resistance to semantic displacement.

\[
\mathcal{K} = g^{ij}\mathrm{Ric}_{ij}.
\]

\begin{proposition}[Semantic Energy Bound]
Any stable RSVP--Polyxan coupling must satisfy
\[
\|\nabla\Phi\|^2 \;\le\; C_1 + C_2 |\mathcal{K}|
\]
for constants $C_1,C_2 > 0$.
\end{proposition}

\begin{proof}[Sketch]
The RSVP Lagrangian includes curvature-weighted entropy terms; unbounded
semantic curvature would destabilize the action functional. See Ch.~2 for full
Lagrangian analysis.
\end{proof}

This bound plays a central role in the global existence of RSVP flows.

% ================================================================
\section{Summary}

We have established:
\begin{itemize}
  \item the semantic manifold $\mathcal{M}$ as the base geometric space,
  \item canonical coordinates reflecting RSVP structure,
  \item jet extensions for higher-order action principles,
  \item stratified semantics and degenerate metrics,
  \item the semantic sheaf $\mathscr{S}$ and semantic fibration $\pi$,
  \item Polyxan atoms as sheaf sections,
  \item semantic flows induced by the vector field $\mathbf{v}$,
  \item curvature bounds ensuring stable RSVP dynamics.
\end{itemize}

These structures prepare the ground for Chapter~2, which introduces the RSVP
action functional, Lagrangian density, and Euler--Lagrange field equations using
the formal apparatus built here.
